% LONG VERSION (kept for other use). Benchmarking and evaluation: forecasting evaluation protocols.
% Draft text for the aeon paper (2026-10-27). References in evaluation_protocols.bib.
% Numbers are from aeon 8585c7bf: results/aeon-8585c7bf/m4/owa_all_forecasters.csv
% and results/aeon-8585c7bf/aaltd/summary.csv in the horizon repository.
% Table labels (tab:m4, tab:aaltd) are placeholders.

\subsection{Forecasting evaluation protocols}
\label{sec:protocols}

Unlike classification and regression, where a train/test split of independent
cases is largely standard, there is no single accepted protocol for evaluating
forecasters. A protocol fixes several choices, each of which can change the
ranking of methods: where the forecast origins are placed (one fixed origin, or
several rolling origins); whether a model is refitted at each origin; how far
ahead it forecasts and how multi-step forecasts are produced; whether a model
is fitted to one series at a time (local) or to a collection of series
(global), possibly after pretraining on other data; and how errors are scaled
and aggregated across series~\citep{tashman2000outofsample,hewamalage2023forecast}.
Multi-step forecasts can be produced \emph{recursively} (also called
\emph{iterated} or \emph{iterative}), by feeding a one-step model's predictions
back in as inputs; \emph{directly}, with a separate model for each step ahead;
or by a single model that outputs the whole horizon at
once~\citep{taieb2012review}. A rolling-origin evaluation, in turn, moves the
forecast origin through a test period and forecasts from each origin, with or
without refitting~\citep{tashman2000outofsample}. These are separate choices
that are often conflated: a rolling evaluation may use recursive, direct or
multi-output forecasts at each origin. aeon exposes the forecasting strategies
explicitly: every forecaster has \texttt{iterative\_forecast} (recursive) and
\texttt{direct\_forecast} methods, and multi-output models
\texttt{series\_to\_series\_forecast}, so the same model can be evaluated
under different protocols.

\paragraph{M4.}
The M4 competition~\citep{makridakis2020m4} is the most widely used forecasting
benchmark. It contains 100{,}000 series at six frequencies, each with a single
fixed origin: the last $h$ observations of every series are held out, with
$h=6, 8, 18, 13, 14$ and $48$ for yearly, quarterly, monthly, weekly, daily and
hourly series. The protocol embeds several assumptions. Each series is
forecast once, from one origin, over the whole horizon, so the score mixes
one-step accuracy with the accumulation of error over up to 48 steps. The
series are short (the minimum training lengths are 13, 16, 42, 80, 93 and 700
observations), and 95\% of them are yearly, quarterly or monthly, so the
overall score is dominated by short, trended, low-frequency series. Accuracy is
summarised by the Overall Weighted Average (OWA), the mean of sMAPE and MASE
each taken relative to Naive2, a seasonal random walk applied to
deseasonalised series, where seasonality is detected by a 90\% autocorrelation
test and removed by classical multiplicative decomposition. The benchmark
methods (including the ML benchmarks, an MLP and an RNN) were applied to series
deseasonalised in this way. OWA is a ratio of means over all series, so
frequencies with few series (weekly and hourly, under 1\% of M4) carry almost
no weight. Finally, the series are not aligned in time: the test period of one
series can lie before the training period of another, which allows information
to leak when global models are trained across series~\citep{ingel2020correlated,hewamalage2023forecast}.
M4 was won by a hybrid of exponential smoothing and a recurrent network trained
across series~\citep{smyl2020hybrid}; purely local machine learning methods
performed poorly~\citep{makridakis2018m4results}.

\paragraph{The Monash repository.}
The Monash Time Series Forecasting Repository~\citep{godahewa2021monash}
(\url{https://forecastingdata.org/}) collects datasets of related series from
many domains, including the M competitions, in a common format, and was
designed for evaluating global models. It keeps M4's single fixed origin: the
last $h$ observations of each series form the test set, with $h$ taken from the
original competition where there is one and otherwise set by frequency (for
example 12 for monthly, 30 for daily and one week, 168, for hourly series).
Baselines are evaluated by MASE, sMAPE, a modified sMAPE robust to zeros, MAE
and RMSE. Many of its series are far longer than those of M4, at high
frequencies with multiple seasonal periods, but each is still evaluated from a
single origin.

\paragraph{M4 in deep learning papers.}
Deep learning papers usually evaluate on M4 under a different protocol.
TimesNet~\citep{wu2023timesnet}, through its code release, the
Time-Series-Library, follows N-BEATS~\citep{oreshkin2020nbeats}: one global
model per frequency is trained on windows sampled from the end of all the
training series of that frequency, it forecasts the whole horizon directly
from only the last $2h$ observations of each series (12 for yearly series, 96
for hourly), and results are reported for yearly, quarterly and monthly series
with weekly, daily and hourly pooled into one ``Others'' group and a weighted
average over all series. The test points and the OWA formula match M4, but the
information each model sees does not: a local statistical model uses the full
history of one series, the deep model a short window of every series of the
same frequency. Moreover, the library's M4 validation set is the M4 test set:
\texttt{Dataset\_M4} loads the test split for any flag other than
\texttt{train}, and the validation loss that drives early stopping and
checkpoint selection scores forecasts from the end of each training series
against the test values. The reported score is therefore that of the
checkpoint selected on the test set. This is present from the library's first
commit (\texttt{e08ddc0}, February 2023), the code release for TimesNet, and
in the current version, so it applies to every model evaluated on M4 with the
library's scripts. It is a model-selection leak, not training on test data,
and its effect on the published scores has not been measured.

\paragraph{GIFT-Eval.}
GIFT-Eval~\citep{aksu2024gifteval} was designed for comparing foundation models
with statistical and deep learning baselines. It combines several protocols.
The final 10\% of each series is used for testing, split into non-overlapping
windows the length of the horizon, so the forecast origin rolls through the
test period in steps of $h$ (at most 20 windows, set by the shortest series in
a dataset); the M4 datasets keep their single original origin and horizons.
Horizons are defined per frequency in three settings, the medium and long
settings being 10 and 15 times the short one (for hourly data, 48, 480 and 720
steps). Within a window each model forecasts the whole horizon in its own way:
local statistical models are fitted to the history before the window (truncated
to the last 1{,}000 observations, with Seasonal Naive substituted when a model
times out); deep learning models are trained once per dataset and condition on
a fixed context; and foundation models forecast zero-shot from a context
window. Decoder-only foundation models such as Chronos and TimesFM produce long
horizons recursively, and the GIFT-Eval authors attribute part of their decline
at long horizons to this~\citep{aksu2024gifteval}, whereas others output the
whole horizon directly. A single GIFT-Eval score therefore aggregates
rolling-origin evaluation, recursive and direct multi-step forecasting, and
local, global and pretrained models, over horizons from one to hundreds of
steps.
% TODO check against the current leaderboard README: metrics (MASE and CRPS
% relative to Seasonal Naive, geometric mean across configurations, ranks by
% CRPS/MASE), as the arXiv v2 text describes MAPE and CRPS.

\paragraph{Our evaluation.}
To show how the protocol changes the conclusions, we evaluate aeon's
forecasters under two contrasting protocols with the same code. On M4 we follow
the competition exactly: one origin per series, the M4 horizons, OWA relative
to Naive2, and M4-style deseasonalisation for the non-seasonal and machine
learning forecasters; the machine learning forecasters predict one step and are
applied recursively. Our second benchmark is a collection of 100 long series
(up to 10{,}000 observations, most from the Monash repository, including 57 long
M4 series) evaluated by rolling origin: each forecaster is refitted at each of
the last 30 time points and makes a one-step forecast, and accuracy is the mean
absolute error relative to the naive forecast, aggregated by geometric mean.
The orderings differ sharply. On M4 the statistical forecasters are best (SCUM
OWA 0.845, AutoETS 0.914) and every machine learning or deep learning forecaster
is worse than Naive2 (DrCIF applied to changes 1.450); on the long series with
one-step rolling-origin forecasts the reverse holds, and DrCIF is the most
accurate (relative MAE 0.781, against 0.843 for SCUM)
(Tables~\ref{tab:m4} and~\ref{tab:aaltd}). Neither result is wrong: they
answer different questions, and a benchmark summarised by one number embeds
the answer to one of them.
